\chapter{Background and Literature Review}
This chapter establishes the theoretical foundation and contextual background for the research presented in this thesis, tracing the technological progression from foundational Optical Transport Networks (OTN) to modern automated control architectures. The review begins by examining the evolution of Wavelength Division Multiplexing (WDM), detailing the industry-wide transition from rigid, fixed-grid spectrum allocation to the highly efficient, elastic paradigms of flexible-grid (flex-grid) EONs. It contrasts traditional, decentralized optical control planes with the emergence of Software-Defined Networking (SDN), which provides the centralized intelligence required for dynamic resource management. Within this context, the chapter details the architectural components of contemporary SDN controllers, specifically focusing on the TeraFlowSDN ecosystem. A critical analysis of current Routing and Spectrum Allocation (RSA) methodologies follows, identifying a persistent research gap: the inability of traditional path computation and spectrum assignment logic to scale effectively when encountering multi-hop topologies with parallel optical links. Finally, the chapter introduces the \texttt{ParallelOpticalController} as a transformative solution within TeraFlowSDN, designed to resolve these computational bottlenecks and provide high-speed, constraint-aware provisioning for complex multifiber infrastructures
\section{Software Defined Network}
\subsection{SDN Architecture and Principles}
Software-Defined Networking (SDN) represents a fundamental paradigm shift in telecommunications, decoupling the network's control logic—the control plane—from the underlying packet or optical forwarding hardware, known as the data plane. This structural separation centralizes network intelligence into a software-based controller that maintains a global, real-time view of the network topology and its available resources. By abstracting the physical infrastructure through open, standardized Southbound Interfaces (SBIs) such as NETCONF or gNMI, SDN enables deep network programmability and dynamic orchestration of forwarding behaviors. This centralized management framework allows high-level applications to request network resources through Northbound Interfaces (NBIs) without requiring granular knowledge of the physical layer's complexity. Consequently, SDN provides the necessary architecture to achieve multi-vendor interoperability and agile service provisioning, transforming rigid physical infrastructure into flexible, automated, and software-governed assets.
\subsection{SDN in Optical Networks}
While initially designed for packet-switched networks, SDN principles have been aggressively adapted for the optical domain to create Software-Defined Optical Networks (SDON). Unlike packet networks, optical networks must manage analog physical layer impairments, wavelength continuity, and complex spectrum allocation constraints. SDON abstracts these physical layer complexities by decoupling the intelligent control logic—termed Optics Defining Software (ODS)—from the programmable physical hardware, known as Software Defined Optics (SDO) \cite{softwareDefinedOpticalNetwork}. This abstraction allows the central controller to dynamically provision lightpaths, tune variable transponders for specific modulation formats, and automatically configure advanced Reconfigurable Optical Add-Drop Multiplexers (ROADMs), successfully moving away from historically static, manually operated optical grids. To manage this complexity at scale, modern SDON deployments frequently adopt a Transport-SDN (T-SDN) architecture, this is often achieved through a hierarchical control structure \cite{largeScaleOpticalSDN}. In this model, lower-level controllers manage specific optical nodes, while a higher-level controller (such as ONOS) facilitates multi-vendor integration using standardized interfaces like the Transport API (T-API) over NETCONF and RESTCONF protocols.
Furthermore, breaking vendor lock-in is a critical focus of modern optical control planes and initiatives like the Open Disaggregated Transport Network (ODTN) utilize standardized YANG data models (such as OpenConfig and OpenROADM) to abstract heterogeneous hardware \cite{odtn}. This standardization enables the seamless, end-to-end integration of packet and optical layers across disaggregated white-box devices. Finally, as emphasized in "Network Slicing in SDN Networks," this centralized, programmable control is what ultimately enables advanced network virtualization. By leveraging SDON capabilities, operators can deploy multiple isolated logical slices over a single shared physical infrastructure, dynamically shaping traffic to guarantee specific Quality of Service (QoS) and performance isolation for varying applications.
\section{TeraflowSDN}
\subsection{Architecture and Ecosystem}
TeraFlowSDN (TFS) is a state-of-the-art, open-source SDN controller developed under the European Telecommunications Standards Institute (ETSI) to radically advance the orchestration of beyond 5G and 6G networks. Designed as a secure, cloud-native platform, TFS employs a microservices-based architecture orchestrated within Kubernetes environments, enabling individual control components to operate autonomously and scale independently. This inherent modularity provides an ideal ecosystem for integrating custom logic, such as vendor-agnostic device drivers and novel path computation engines. Furthermore, the architecture is explicitly engineered to seamlessly integrate with modern Network Function Virtualization (NFV) and Multi-access Edge Computing (MEC) frameworks, including ETSI OpenSourceMANO. By unifying network and cloud resource management, TeraFlowSDN equips telecom operators and hyperscale cloud providers with zero-touch automation, high business agility, and dynamic service provisioning across complex, multi-tenant optical and packet transport infrastructures.
\begin{figure}[htbp]
    \centering
    \includegraphics[width=0.6\textwidth]{tfs_architecture.png}
    \caption{TeraflowSDN Architecture}
    \label{fig:tfs_architecture}
\end{figure}
\subsection{Key services}
The TeraFlowSDN ecosystem relies on a highly coordinated suite of specialized microservices, communicating via Remote Procedure Calls (gRPC), to maintain network state and execute dynamic configurations. At the core of this architecture is the Context Service, which acts as the highly available, stateful storage mechanism maintaining the definitive global view of multi-domain topologies, devices, endpoints, and active network slices. Interacting closely with this centralized database is the Device Service, which manages all southbound interfaces by utilizing a modular framework of vendor-agnostic drivers—often leveraging NETCONF and OpenConfig models—to dynamically onboard and configure physical or emulated hardware. Furthermore, the Service Component oversees the end-to-end lifecycle of network connections by interpreting incoming connectivity requests and delegating cross-layer configurations to appropriate handlers. For optical transport infrastructures specifically, the OpticalController operates as a dedicated service to process complex optical intents, executing resource allocation algorithms such as Routing and Spectrum Assignment (RSA) to establish continuous and contiguous lightpaths across multi-hop networks. These fundamental operations are visually managed through a graphical WebUI that provides topology visualization and telemetry display, while the broader ecosystem is fortified by complementary microservices dedicated to machine learning-based security, distributed ledger technologies for multi-tenancy, and policy-driven automation to ensure robust, zero-touch network orchestration.
\section{Optical Transport Networks}
\subsection{WDM technology}
Wavelength Division Multiplexing (WDM), specifically Dense WDM (DWDM), is the foundational technology for modern high-capacity optical networks. DWDM enables the transmission of multiple data streams over a single optical fiber by modulating each stream onto a distinct wavelength (or color) of laser light. Historically, DWDM utilized a fixed-grid spacing, allocating rigid channel intervals of 50 GHz, 100 GHz, or even wider. However, as data rates increased and traffic patterns became more unpredictable, this rigid spacing resulted in inefficient spectrum utilization. To address this bottleneck, modern elastic optical networks (EONs) utilize flexible-grid (flex-grid) technologies\cite{opticalNetworksAndMangementControl}. As defined in the ITU-T G.694.1 standard, the flex-grid allows a mixed bit rate or mixed modulation format transmission system to allocate frequency slots with different widths, optimizing the spectrum for the specific bandwidth requirements of individual channels\cite{itu694}. This transition from fixed to flexible allocation is what necessitates advanced Routing and Spectrum Allocation (RSA) algorithms, such as the one proposed in this thesis, to dynamically manage the resulting variable-width spectrum blocks across complex, multi-fiber topologies.
\begin{figure}[htbp]
    \centering
    \includegraphics[width=0.6\textwidth]{dwdm_flex_grid.jpg}
    \caption{DWDM Flex-Grid Overview}
    \label{fig:dwdm_flex_grid}
\end{figure}
\subsection{ITU standards}
The implementation of these flex-grid networks is strictly governed by International Telecommunication Union (ITU) standards to ensure interoperability and consistent spectrum management. Specifically, the ITU-T G.694.1 Recommendation provides the definitive frequency grid for DWDM applications. Under this standard, the entire DWDM spectrum is anchored to a reference central frequency of 193.1 THz.
To enable elastic allocation, the standard defines a "frequency slot" as a contiguous frequency range characterized by its nominal central frequency and its total slot width. The architecture is built upon two critical granularities:
\begin{itemize}
    \item \textbf{Nominal Central Frequency Granularity:} The allowed central frequencies are defined by the formula $193.1 + n \times 0.00625$ THz, establishing a fine-grained granularity of 6.25 GHz to allow precise positioning of channels.
    \item \textbf{Slot Width Granularity:} The total width of any assigned frequency slot is defined by $12.5 \times m$ GHz, establishing 12.5 GHz as the minimum, indivisible building block for bandwidth allocation.
\end{itemize}
\subsection{Network elements (devices)}
The physical topology comprises two primary categories of optical nodes: terminal multiplexers (transponders and muxponders) and Reconfigurable Optical Add-Drop Multiplexers (ROADMs). In modern high-capacity optical networks, network endpoints are more precisely characterized as muxponders, which aggregate multiple lower-rate client-layer signals into high-speed optical channels suitable for DWDM transmission. These devices utilize bidirectional ports equipped with advanced pluggable transceivers that handle complex modulation and Forward Error Correction (FEC). While the telecommunications industry employs a diverse array of transceiver form factors and optical modules—such as QSFP, QSFP28, OpenZR+, and 400ZR—each offering varying frequency band support, capacity, and error correction capabilities, this thesis adopts a specific, standardized hardware profile for its emulation environment. Specifically, the emulated muxponder ports are modeled using the QSFP56-DD form factor equipped with Digital Coherent Optics, utilizing a 400GBASE-ZR Ethernet Physical Medium Dependent (PMD) sublayer and automated FEC. Although real-world deployments may feature heterogeneous transceiver capabilities, a foundational assumption of this research is that all muxponder ports are fully bidirectional and provide continuous, tunable frequency support across the entire C-band. This baseline ensures a uniform spectrum allocation space for evaluating the proposed RSA algorithms.
\begin{figure}[htbp]
    \centering
    \includegraphics[width=0.6\textwidth]{roadm_muxponder.png}
    \caption{An example optical network}
    \label{fig:roadm_muxponder}
\end{figure}
Conversely, ROADMs serve as the optical switching backbone, enabling individual wavelengths to be dynamically added, dropped, or passed through a node entirely in the optical domain, thereby bypassing costly optical-to-electrical-to-optical (O-E-O) conversions. To accurately evaluate the Routing and Spectrum Allocation (RSA) algorithm on parallel links, it is necessary to establish strict assumptions regarding the internal architecture of these ROADMs. A typical ROADM architecture consists of Shared Risk Groups (SRGs) for local add/drop functionality and Direction Degrees (DEGs) equipped with Wavelength Selective Switches (WSS) for line-side routing. In the context of this thesis, physical fiber connections are only classified as strict "parallel links" if the ports originating from the same muxponder terminate at the same SRG and are routed outward through the exact same DEG of the ROADM. Under this specific topological constraint, the parallel fibers share an identical physical routing vector, making spectrum contiguity and continuity algorithms highly complex, as the spatial routing does not inherently resolve spectrum overlap. If the ports were to route through different DEGs, the spatial separation provided by distinct WSS hardware would allow for simple frequency re-use, thereby negating the specific parallel-link RSA constraints addressed in this study. Furthermore, to accommodate high-capacity throughput and scalability, the ROADM devices in this simulated environment are assumed to be wideband-capable, operating seamlessly across the entire S, C, and L (SCL) optical bands.
\subsection{Data models (configuration)}
To achieve true interoperability and break the traditional hardware lock-in imposed by proprietary management interfaces, modern optical equipment relies on standardized YANG data models. OpenConfig\cite{openconfigYang} and OpenROADM\cite{openroadmYang} have emerged as the two prominent industry standards used to formalize the representation of optical device configurations and operational states. By structuring parameters in universally understood XML or JSON formats, these models hide the underlying hardware idiosyncrasies from the control plane, allowing an SDN controller to parse telemetry and push configurations uniformly regardless of the equipment manufacturer. To manage the complexity of optical networking, these data models utilize strictly defined XML namespaces to isolate specific operational domains. Within the emulation environment of this thesis, the terminal devices (muxponders) primarily utilize the OpenConfig standard. For instance, the physical hardware hierarchy—such as form factors and transceivers—is defined under the \href{http://openconfig.net/yang/platform}{plaform-namespace}. Concurrently, the operational optics are managed via the \href{http://openconfig.net/yang/terminal-device}{terminal-device namespace}, which dictates precise parameters such as operational frequencies (e.g., 191.6 THz), target output power, and the logical channel assignments mapping Ethernet client traffic to Optical Transport Network (OTN) line ports.

Conversely, the emulated optical switches utilize the OpenROADM data model, which is specifically tailored to describe complex photonic switching architectures. Using namespaces such as \href{http://org/openroadm/device}{device namespace} and \href{http://org/openroadm/interfaces}{interface namespace}, the model accurately maps the internal ROADM topology. It defines critical switching components including circuit packs (such as twin Wavelength Selective Switches), Direction Degrees (DEGs), Shared Risk Groups (SRGs), and specific optical interfaces like Optical Transport Section (OTS) and Optical Multiplex Section (OMS). By adhering to these standardized schemas, the SDN controller is decoupled from vendor-specific command-line interfaces. Consequently, network operators can seamlessly deploy heterogeneous, multi-vendor optical topologies. A single controller can dynamically orchestrate an end-to-end lightpath by pushing standard OpenConfig XML payloads to tune a Cisco transceiver, while simultaneously pushing OpenROADM payloads to configure cross-connects on an Ericsson ROADM, effectively unifying the entire physical layer.
\section{Routing and Spectrum Allocation}
\subsection{Spectrum continuity and contiguity constraints}
Establishing a lightpath in an Elastic Optical Network (EON) requires solving the complex Routing and Spectrum Allocation (RSA) problem. In traditional fixed-grid DWDM networks, resource allocation—known as Routing and Wavelength Assignment (RWA)—was primarily concerned with satisfying a simple wavelength continuity requirement \cite{rsaHeuristicSliced}. However, the evolution to flexible-grid EONs introduces highly restrictive dual constraints that transform resource allocation into an NP-hard computational challenge \cite{rsaHeuristicSliced}. Regardless of the specific allocation policy utilized, every valid assignment in an EON must fundamentally satisfy three primary physical constraints \cite{algorithmsRSATaxonomy}:

\begin{itemize}
    \item \textbf{Spectrum Continuity:} This constraint dictates that the exact same frequency spectrum (the specific index of the allocated spectrum slots) must be utilized and remain identical across every single link along the end-to-end routing path \cite{algorithmResearchRSAONR}. However, in multifiber infrastructures featuring parallel links, this constraint presents a crucial operational nuance \cite{efficientDynamicRSAEon}. While the frequency index must remain identical end-to-end, the optical signal is permitted to shift spatially between parallel physical fibers at each hop, provided there are no internal switching restrictions at the optical cross-connects \cite{efficientDynamicRSAEon}.
    \item \textbf{Spectrum Contiguity:} The frequency slots assigned to a single lightpath must be strictly adjacent to one another within the frequency domain to form a contiguous "Slot Block". In modern implementations, algorithms treat individual frequency slots as fixed at a minimum width (typically 12.5 GHz) and group them contiguously to precisely match the variable bandwidth demands of the traffic \cite{algorithmsRSATaxonomy}.
    \item \textbf{Spectrum Non-Overlapping:} To prevent signal interference and data corruption, this constraint ensures that any specific frequency slot on a given physical fiber is allocated to a maximum of one active lightpath at a time \cite{efficientDynamicRSAEon}.
\end{itemize}
The tight coupling of the continuity and contiguity constraints often leads to severe spectrum fragmentation. Because of the strict contiguity requirement, having enough total available spectrum on a path is insufficient if those slots are scattered and fragmented. To enforce these rules without blocking new requests, advanced RSA algorithms manage both constraints simultaneously by treating the network state as a unified "allocation matrix". This allows the controller to search for the first available block of contiguous slots that spans uniformly across all physical edges along the chosen multi-hop route \cite{rsaFlexGrid}.
\subsection{Existing RSA algorithms}
As established, the strict enforcement of spectrum continuity and contiguity transforms network resource allocation into a highly restrictive, NP-hard computational challenge\cite{rsaHeuristicSliced}. To prevent severe spectrum fragmentation and the unfair blocking of high-capacity requests, intense research has been dedicated to developing sophisticated Routing and Spectrum Allocation (RSA) algorithms\cite{distanceAdaptive}. Over time, these algorithms have evolved from basic wavelength assignment models into complex, multi-metric optimization strategies designed to maximize network throughput and accommodate diverse, dynamic traffic demands\cite{algorithmsRSATaxonomy}. Early approaches sought to decouple the routing and assignment phases to mitigate computational bottlenecks. Foundational algorithms employed greedy heuristics and K-shortest paths to maximize single-hop traffic while strictly enforcing standard wavelength continuity\cite{heuristicWSAMultihop}. As networks transitioned to flex-grid EONs, researchers introduced offline traffic sorting heuristics to groom traffic more efficiently. For instance, the Minimum Hops with Least Slots (MHLS) algorithm prioritizes connection requests by calculating a sorting parameter\cite{rsaHeuristicSliced}:
\[
    NP=Hops \times Slots
\]
By serving "smaller" requests first, the algorithm minimizes overall blocked traffic\cite{rsaHeuristicSliced}. More dynamic approaches have integrated custom cost functions to balance multiple network metrics simultaneously. For example, the \texttt{FRSO-RSA} algorithm evaluates candidate alternative routes by actively steering traffic away from congested or unreliable paths using the following cost function\cite{failureProbability}:
\[
    cost_c = \alpha \times PF_c + (1 - \alpha) OU_c
\]
where $PF_c$ represents the physical failure rate of the path, $OU_c$ denotes the spectrum occupancy rate, and $\alpha$ acts as an adjustable weight parameter\cite{failureProbability}.
To combat the specific issue of contiguity-induced fragmentation, recent literature has shifted toward "next-state-aware" algorithms. Metrics such as Maximum Spectrum Completeness (MSC) are used to evaluate how a current assignment will impact future capacity, actively selecting routes that preserve large, unbroken blocks of spectrum \cite{maximumSpectrumCompleteness}. Other strategies introduce highly dynamic link-weighting mechanisms that actively draw traffic toward links harboring the largest unfragmented spectrum slices rather than just looking at total available capacity \cite{spectrumStatusAware}. When fragmentation inevitably occurs during the network's lifecycle, complex defragmentation algorithms utilizing Integer Linear Programming (ILP) and Resource Dependency Digraphs (RDD) are deployed to mathematically recalculate and shift existing lightpaths without altering their physical routes\cite{optimizingSpectrumAllocationEON}. Furthermore, advanced meta-heuristics like Genetic Algorithms and probabilistic models like Markov Decision Processes (MDP) have been explored to navigate massive solution spaces and predict network states\cite{routingModulation,adaptiveRSAFlex}.
\section{Problems with Parallel Links in RSA}
\subsection{Problem statement}
While the Routing and Spectrum Allocation (RSA) problem in optical networks has been extensively studied for years, the comprehensive integration of multi-fiber or parallel links between adjacent nodes remains an underexplored area. Technologies such as link aggregation and bandwidth duplexing have become increasingly common, introducing multiple physical connections between devices to meet escalating capacity demands. Specifically, to aggregate a higher number of optical channels, multiple parallel connections are routinely established between a muxponder and a Reconfigurable Optical Add-Drop Multiplexer (ROADM). A critical physical constraint in this architecture is that the same spectrum cannot be reused across parallel links that terminate at the same degree of a ROADM device without causing severe interference. Traditional RSA algorithms primarily focus on path computation, spectrum continuity, and spectrum contiguity across single-fiber logical links. However, they frequently fail to account for the spatial complexity introduced by multi-fiber or parallel-link topologies. Current RSA algorithms already face significant scalability and efficiency drawbacks in complex, highly meshed topologies; introducing parallel links exacerbates these challenges dramatically. In a multi-graph network featuring a uniform or variable number of parallel links between adjacent devices, the total number of possible path combinations from a source to a destination increases exponentially. Because signals routed over the same ROADM degree cannot share frequencies, a network controller must exhaustively calculate and verify the availability of every frequency slot across every parallel edge. Computing paths in such multi-dimensional scenarios can easily overwhelm current algorithms, resulting in massive computational overhead, unacceptable latency, and ultimately, a failure in dynamic lightpath establishment. Consequently, there remains a critical research gap in developing routing strategies that can navigate the spatial complexity of parallel links without crippling network efficiency.
\section{Current approaches and limitation}
To address the NP-hard complexities of optical network management, theoretical research has developed sophisticated Routing and Spectrum Allocation (RSA) algorithms designed to mitigate severe spectrum fragmentation and maximize throughput. These strategies have evolved from foundational, decoupled wavelength assignment models and offline traffic sorting heuristics, into dynamic, multi-metric optimization approaches that actively steer traffic away from congested or unreliable paths. Recent literature also heavily emphasizes "next-state-aware" metrics, dynamic link-weighting, and complex defragmentation models. However, extensive taxonomic reviews reveal a prevailing structural limitation: the vast majority of these theoretical solutions are designed exclusively for one-dimensional, single-fiber-per-link topologies. When applied to multi-fiber scenarios, these models must treat parallel links either as physically invalid aggregated bandwidth or as entirely separate logical links, inevitably triggering a state-space explosion and unmanageable computational complexity.

In practical deployments, industry-standard SDN controllers like ONOS\cite{odtnONOS} and OpenDaylight (ODL)\cite{odlTPCE} attempt to manage these topologies using robust, graph-based routing engines, yet they fundamentally struggle with the same spatial constraints. In both controllers, optical connections are treated as strictly matched port-to-port links; thus, multiple parallel fibers between the same devices are mapped as entirely distinct and independent topological edges. When computing routes via standard K-Shortest Path algorithms, the topology engine outputs every permutation of these parallel links as a completely independent path sequence. Consequently, the subsequent spectrum allocation phase applies strict isolation, evaluating spectrum continuity and contiguity sequentially on a strict per-path basis. If an allocation attempt experiences wavelength contention on one parallel link, the controller discards the effort and must iteratively re-run the entire RSA verification from scratch on the next topological path permutation. In complex, highly meshed topologies featuring multiple parallel links, this exhaustive, sequential path-by-path iteration causes exponential computational overhead, often leading to unacceptable latency and dynamic provisioning timeouts. These systemic shortcomings in both theoretical models and practical SDN implementations highlight a critical research gap, presenting the opportunity to design a deterministic, spatially-aware parallel optical controller within the TeraFlowSDN ecosystem that evaluates parallel ports concurrently to bypass sequential iteration and enable highly scalable lightpath establishment.
\section{Related work}
While the Routing and Spectrum Allocation (RSA) problem has been extensively explored to address strict spectrum continuity and contiguity constraints, existing literature and practical deployments exhibit a critical blind spot regarding multi-fiber topologies. Advanced theoretical heuristics—such as next-state-aware metrics and dynamic link-weighting—are overwhelmingly optimized for one-dimensional logical links. When extended to parallel-link environments, both theoretical models and industry-standard SDN controllers typically resort to exhaustive, sequential path-by-path iterations or restrictive offline computations, leading to severe latency and state-space explosion. To bridge this gap, this thesis introduces the \texttt{ParallelOpticalController} within the TeraFlowSDN ecosystem. By concurrently evaluating all available parallel ports at each hop rather than processing them as isolated topological sequences, the proposed architecture natively leverages the spatial domain to resolve RSA constraints in real time, bypassing sequential bottlenecks and enabling highly scalable, low-latency lightpath establishment for next-generation multi-fiber networks.